\documentclass[10pt,a4paper]{amsart}
\usepackage{amsmath,amssymb,amsthm,graphicx}
% Retained minimal preamble for the Stacks Project chapter more-algebra.tex (More on Algebra).
% The chapter file pinned-original-source.tex carries only \input{preamble} on its first line, so no
% preamble text lives inside the chapter itself.  The Stacks Project uses one single shared preamble for
% every chapter; the snapshot stages it identically as sources/stacks-chow/preamble.tex and
% sources/stacks-etale-cohomology/preamble.tex.  The block below is copied mechanically, line for line,
% from sources/stacks-chow/preamble.tex lines 14-18 (the Xy-pic setup) and 154-217 (the native theorem
% declarations together with the mathematical macros).  Nothing was added, removed or redefined; the
% wrapper class and the amsmath/amssymb/amsthm/graphicx base packages are supplied by the assembler.
% Chapter: more-algebra.tex -- More on Algebra.
% For commutative diagrams we use Xy-pic
\usepackage[all]{xy}

% We use 2cell for 2-commutative diagrams.
\xyoption{2cell}

% Theorem environments.
%
\theoremstyle{plain}
\newtheorem{theorem}[subsection]{Theorem}
\newtheorem{proposition}[subsection]{Proposition}
\newtheorem{lemma}[subsection]{Lemma}

\theoremstyle{definition}
\newtheorem{definition}[subsection]{Definition}
\newtheorem{example}[subsection]{Example}
\newtheorem{exercise}[subsection]{Exercise}
\newtheorem{situation}[subsection]{Situation}

\theoremstyle{remark}
\newtheorem{remark}[subsection]{Remark}
\newtheorem{remarks}[subsection]{Remarks}

\numberwithin{equation}{subsection}

% Macros
%
\def\lim{\mathop{\mathrm{lim}}\nolimits}
\def\colim{\mathop{\mathrm{colim}}\nolimits}
\def\Spec{\mathop{\mathrm{Spec}}}
\def\Hom{\mathop{\mathrm{Hom}}\nolimits}
\def\Ext{\mathop{\mathrm{Ext}}\nolimits}
\def\SheafHom{\mathop{\mathcal{H}\!\mathit{om}}\nolimits}
\def\SheafExt{\mathop{\mathcal{E}\!\mathit{xt}}\nolimits}
\def\Sch{\mathit{Sch}}
\def\Mor{\mathop{\mathrm{Mor}}\nolimits}
\def\Ob{\mathop{\mathrm{Ob}}\nolimits}
\def\Sh{\mathop{\mathit{Sh}}\nolimits}
\def\NL{\mathop{N\!L}\nolimits}
\def\CH{\mathop{\mathrm{CH}}\nolimits}
\def\proetale{{pro\text{-}\acute{e}tale}}
\def\etale{{\acute{e}tale}}
\def\QCoh{\mathit{QCoh}}
\def\Ker{\mathop{\mathrm{Ker}}}
\def\Im{\mathop{\mathrm{Im}}}
\def\Coker{\mathop{\mathrm{Coker}}}
\def\Coim{\mathop{\mathrm{Coim}}}

% Boxtimes
%
\DeclareMathSymbol{\boxtimes}{\mathbin}{AMSa}{"02}

%
% Macros for moduli stacks/spaces
%
\def\QCohstack{\mathcal{QC}\!\mathit{oh}}
\def\Cohstack{\mathcal{C}\!\mathit{oh}}
\def\Spacesstack{\mathcal{S}\!\mathit{paces}}
\def\Quotfunctor{\mathrm{Quot}}
\def\Hilbfunctor{\mathrm{Hilb}}
\def\Curvesstack{\mathcal{C}\!\mathit{urves}}
\def\Polarizedstack{\mathcal{P}\!\mathit{olarized}}
\def\Complexesstack{\mathcal{C}\!\mathit{omplexes}}
% \Pic is the operator that assigns to X its picard group, usage \Pic(X)
% \Picardstack_{X/B} denotes the Picard stack of X over B
% \Picardfunctor_{X/B} denotes the Picard functor of X over B
\def\Pic{\mathop{\mathrm{Pic}}\nolimits}
\def\Picardstack{\mathcal{P}\!\mathit{ic}}
\def\Picardfunctor{\mathrm{Pic}}
\def\Deformationcategory{\mathcal{D}\!\mathit{ef}}
\title{Vanishing of H\textasciicircum{}0 of the product totalization of a complex of complexes whose p-th term has vanishing H\textasciicircum{}\{-p\}}
\author{The Stacks Project Authors}
\date{stacks-project:more-algebra.tex:local-snapshot-20261012; GNU FDL 1.2 or later}
\begin{document}
\maketitle
\setcounter{subsection}{1}
\noindent\textit{Editorial note.} All delivered mathematical text is the author's own, copied from the pinned Stacks Project chapter file more-algebra.tex; no statement or proof step was rewritten, reordered or removed. The counted statement is the labeled lemma environment at source lines 30095-30102 and the counted proof is the proof environment at source lines 30104-30128, which immediately follows it and is its own proof.\par\medskip
\noindent\textit{Editorial note.} The section lead-in at source lines 30092-30093, pointing to the discussion of double complexes in the Homology chapter, lies directly before the counted statement and is outside the counted spans; the card therefore needs no navigation stub, and the proof refers to no labelled item of any chapter.\par\medskip
\noindent\textit{Editorial note.} The chapter file pinned-original-source.tex carries only an input of the shared preamble on its first line, so no preamble text lives inside the chapter. The retained file sources/stacks-more-algebra/retained-preamble.tex is a mechanical, line-for-line copy of the project's single shared preamble as staged in the snapshot (Xy-pic setup plus the native theorem declarations and mathematical macros). Nothing was added, removed or redefined, and the excerpt serves only as preamble, never as a span coordinate.\par\medskip
\noindent\textit{Editorial note.} The source URL is the source-level Stacks Project URL, as required when more than one card is produced from one source; the per-claim tag is carried by claim\_key. The counted statement and proof cite no bibliography entry, so this card needs no bibliography.\par\medskip
\noindent\textit{Allowed standard prerequisites (editorial).}\begin{itemize}
\item The differential on the product totalization Tot\_pi of a double complex with entries A\textasciicircum{}\{p,q\} is, on elements of total degree n, the product over p + q = n of the components f\textasciicircum{}q\_p + (-1)\textasciicircum{}p d\textasciicircum{}q\_\{A\_p\}. recalled at proof lines 30106-30111. Gives the explicit form of the single equation d(x) = 0 for a cocycle representative x = (x\_p) in the product, which is then unwound degree by degree.
\item A class in H\textasciicircum{}0 is represented by a cocycle of total degree zero, i.e. by a family x = (x\_p) in the product of the A\textasciicircum{}\{p,-p\} with d(x) = 0. used at proof lines 30112-30115. Fixes the shape of the object that has to be exhibited as a boundary, so that the goal becomes finding y with d(y) = x.
\item Exactness of a complex at a degree: if H\textasciicircum{}q(C) = 0 and an element z has d(z) = 0, then z is the boundary of some element of C\textasciicircum{}\{q-1\}. used at proof lines 30116-30121, first with H\textasciicircum{}0(A\_0) = 0 and then with H\textasciicircum{}\{-1\}(A\_1) = 0. Supplies the first element y\_\{-1\} with d(y\_\{-1\}) = x\_0 and then the element y\_\{-2\} that kills the corrected component x\_1 + f\_0(y\_\{-1\}); the same step is iterated in the induction.
\end{itemize}\medskip
\begin{lemma}
\label{lemma-vanishing-coh-prod-totalization}
Let $A_0^\bullet \to A_1^\bullet \to A_2^\bullet \to \ldots$
be a complex of complexes of abelian groups. Assume
$H^{-p}(A_p^\bullet) = 0$ for all $p \geq 0$.
Set $A^{p, q} = A_p^q$ and view $A^{\bullet, \bullet}$ as a double complex.
Then $H^0(\text{Tot}_\pi(A^{\bullet, \bullet})) = 0$.
\end{lemma}

\begin{proof}
Denote $f_p : A_p^\bullet \to A_{p + 1}^\bullet$ the given maps
of complexes. Recall that the differential on
$\text{Tot}_\pi(A^{\bullet, \bullet})$ is given by
$$
\prod\nolimits_{p + q = n} (f^q_p + (-1)^p\text{d}^q_{A_p^\bullet})
$$
on elements in degree $n$.
Let $\xi \in H^0(\text{Tot}_\pi(A^{\bullet, \bullet}))$ be a cohomology
class. We will show $\xi$ is zero. Represent $\xi$ as the class of an cocycle
$x = (x_p) \in \prod A^{p, -p}$.
Since $\text{d}(x) = 0$ we find that
$\text{d}_{A_0^\bullet}(x_0) = 0$.
Since $H^0(A_0^\bullet) = 0$ there exists a $y_{-1} \in A^{0, -1}$ with
$\text{d}_{A_0^\bullet}(y_{-1}) = x_0$.
Then we see that $\text{d}_{A_1^\bullet}(x_1 + f_0(y_{-1})) = 0$.
Since $H^{-1}(A_1^\bullet) = 0$ we can find a $y_{-2} \in A^{1, -2}$
such that $-\text{d}_{A_1^\bullet}(y_{-2}) = x_1 + f_0(y_{-1})$.
By induction we can find
$y_{-p - 1} \in A^{p, -p - 1}$ such that
$$
(-1)^p\text{d}_{A_p^\bullet}(y_{-p - 1}) = x_p + f_{p - 1}(y_{-p})
$$
This implies that $\text{d}(y) = x$ where $y = (y_{-p - 1})$.
\end{proof}

\end{document}
